% SGA 3，第 VII_B 篇，英文重建稿——第 1 轮。
% 范围：权威底本分篇第 55 页，第 2.4.2 节全文，止于第 2.5 节之前。
% 合并阅读本页码： 568.
% 再版印刷页码： 557.
% 本范围内原版页边定位： 522.
% 本范围内编者注： 94.
% 权威底本： Polo--Gille Exp7B-13oct24.pdf,
% SHA-256 D303690E0195654BBE7C3DAA5D89EA23E24CD709A3575217234AC19AE8836694.
% 曾参考 Jacob Reinhold 的 CC BY 4.0 Markdown 作为比对谱系。
% 本范围内没有图表。

\subsubsection*{2.4.2}
\phantomsection
\label{sga3:VIIB:2.4.2}

沿用前面的记号，并假设 \(H\) 和 \(G\) 在 \(k\) 上拓扑平坦。则 \(G\) 在 \(k\) 和 \(G/H\) 上均拓扑平坦；故由 \SGARefLink{sga3:VIIB:1.3.3}{1.3.3}，\(G/H\) 在 \(k\) 上拓扑平坦。因此由 \SGARefLink{sga3:VIIB:2.4}{2.4}，典范投影 \[
  q:K\longrightarrow\operatorname{Coker}(f')
\] 拓扑平坦，且 \(f'\) 为从 \(G/H\) 到 \(\ker(q)\) 的同构。又显然有 \[
  \operatorname{Coker}(f)
  =
  \operatorname{Coker}(f').
\]。所以在 \(G\) 和 \(H=\ker(f)\) 于 \(k\) 上拓扑平坦的假设下，已得到 \(f\) 的像 \(\ker(q)\) 与 \(f\) 的余像 \[
  G/\ker(f),
\] 之间的同构。这就证明了下面的定理。

\medskip\noindent\textbf{定理。}\phantomsection \label{sga3:VIIB:2.4.2:theorem}{\itshape 设 \(k\) 为域。交换形式 \(k\)-群构成阿贝尔范畴。}

\medskip\noindent\textbf{推论。}\phantomsection \label{sga3:VIIB:2.4.2:corollary}{\itshape 设 \(k\) 为域。交换仿射 \(k\)-群概形构成阿贝尔范畴。}\footnote{\label{sga3:VIIB:nde:94}另见第 VI\textsubscript{A} 讲，\SGARefLink{cor:sga3-VIA-5.4.3}{5.4.3} 后的注。}

这由上述定理及 \SGARefLink{sga3:VIIB:2.2.2}{2.2.2} 中的范畴等价得出。
